\ReportSection{2. Story Testing Method}
\ReportSubsection{2.1 User Stories and Acceptance Criteria}
We chose a student organizing coursework and personal tasks as our user. Each
story describes something the student wants to do and why it is useful. For
example, Story 6 states: ``As a student, I want to create a category so that I can
organize my work into meaningful groups.'' The acceptance tests then check
whether the API supports that goal.

We used our Part A findings to decide what to check after each operation.
Relationship tests read the association from both objects, deletion tests check
whether references remain, and write tests read back the stored records. This
matters because a request can return a successful status while leaving the data
in an unexpected state. We compared the results with the API documentation and
the expected behavior recorded in our bug report~[\ref{ref:api},\ref{ref:bugs}].

\ReportSubsection{2.2 Normal, Alternate, and Error Flows}
Every feature has a \textit{Background} and at least one normal, alternate, and error
\textit{Scenario Outline}. The \textit{Background} starts the application and creates the data
needed for each example. The normal flow completes the main operation. The
alternate flow reaches the goal with a different valid input or approach. The
error flow checks what happens when the operation cannot be completed.

For Story 6, the normal flow creates a category with a title and description.
The alternate flow creates one with only a title. The error flow submits an empty
title and expects HTTP 400 with the validation message. It also compares the
stored records before and after the request, because rejecting the input should
leave the system unchanged.

Each \textit{Examples} table supplies the data for its \textit{Scenario Outline}. Most outlines
have two rows, while two todo outlines have a third row. Our 72 outlines therefore
produce 146 cases, including the relationship and filter checks that reproduce
our reported findings.

\ReportSubsection{2.3 Deciding What the API Should Do}
Some expectations follow from the documentation or from the operation itself.
For example, asking for a missing project's tasks should not return work from
another project. Our test accepts either HTTP 404 with the expected error shape
or HTTP 200 with an empty task array, because the documentation does not specify
an exact missing-project response for this query.

Other expectations come from how we want the student workflow to behave. These
include seeing category membership from both objects, accepting a numeric ID
where its string form works, and preserving links when changing only a title.
Where the API does not clearly promise this behavior, we mark the failing
scenario with both \Path{@bug} and \Path{@policy}. This way, the result shows
which expectation failed and whether that expectation is an established API rule
or a rule we chose for the workflow. We treat supported inline creation through
a relationship request as a valid alternate flow.
